\documentclass[12pt, a4paper]{article}

% ── PACKAGES ──────────────────────────────────────────────────────────────────
\usepackage[T1]{fontenc}
\usepackage[utf8]{inputenc}
\usepackage[english]{babel}
\usepackage[top=2.5cm, bottom=2.5cm, left=2.8cm, right=2.8cm]{geometry}
\usepackage{lmodern}
\usepackage{microtype}
\usepackage{setspace}
\usepackage{parskip}

% Colours
\usepackage[dvipsnames, table]{xcolor}
\definecolor{forestgreen}{HTML}{2d6a4f}
\definecolor{lightgreen}{HTML}{d8f3dc}
\definecolor{midgreen}{HTML}{52b788}
\definecolor{darkgreen}{HTML}{1a3a28}
\definecolor{warnbg}{HTML}{fff3e0}
\definecolor{warnborder}{HTML}{f57c00}
\definecolor{bluebg}{HTML}{e3f2fd}
\definecolor{blueborder}{HTML}{1565c0}
\definecolor{rowshade}{HTML}{f0ede4}

% Tables
\usepackage{booktabs}
\usepackage{tabularx}
\usepackage{multirow}
\usepackage{array}
\usepackage{longtable}

% Boxes / callouts
\usepackage[most]{tcolorbox}
\tcbuselibrary{breakable, skins}

% Headers / sections
\usepackage{titlesec}
\usepackage{titling}
\usepackage{fancyhdr}

% Lists
\usepackage{enumitem}

% Links
\usepackage[hidelinks, colorlinks=true, linkcolor=forestgreen,
            urlcolor=forestgreen, citecolor=forestgreen]{hyperref}

% Bibliography
\usepackage[numbers, sort&compress]{natbib}

% ── SECTION FORMATTING ────────────────────────────────────────────────────────
\titleformat{\section}
  {\Large\bfseries\color{forestgreen}}
  {\thesection.}{0.8em}{}[{\color{midgreen}\titlerule[1.5pt]}]

\titleformat{\subsection}
  {\large\bfseries\color{darkgreen}}
  {\thesubsection.}{0.7em}{}

\titleformat{\subsubsection}
  {\normalsize\bfseries\color{darkgreen}}
  {\thesubsubsection.}{0.6em}{}

\titlespacing{\section}{0pt}{24pt}{10pt}
\titlespacing{\subsection}{0pt}{16pt}{6pt}

% ── HEADER / FOOTER ───────────────────────────────────────────────────────────
\pagestyle{fancy}
\fancyhf{}
\rhead{\textcolor{forestgreen}{\small Soundscapes \& Biodiversity}}
\lhead{\textcolor{gray}{\small\leftmark}}
\cfoot{\thepage}
\renewcommand{\headrulewidth}{0.4pt}
\renewcommand{\headrule}{\color{midgreen}\hrule width\headwidth height\headrulewidth}

% ── CALLOUT BOXES ─────────────────────────────────────────────────────────────
\newtcolorbox{greenbox}[1][]{
  breakable,
  enhanced,
  colback=lightgreen,
  colframe=midgreen,
  boxrule=0pt,
  leftrule=4pt,
  arc=3pt,
  fonttitle=\bfseries\color{darkgreen},
  coltitle=darkgreen,
  #1
}

\newtcolorbox{warnbox}[1][]{
  breakable,
  enhanced,
  colback=warnbg,
  colframe=warnborder,
  boxrule=0pt,
  leftrule=4pt,
  arc=3pt,
  #1
}

\newtcolorbox{bluebox}[1][]{
  breakable,
  enhanced,
  colback=bluebg,
  colframe=blueborder,
  boxrule=0pt,
  leftrule=4pt,
  arc=3pt,
  #1
}

\newtcolorbox{speciesbox}[2][]{
  breakable,
  enhanced,
  colback=white,
  colframe=forestgreen,
  boxrule=0.8pt,
  arc=4pt,
  title={\bfseries #2},
  fonttitle=\color{white},
  attach boxed title to top left={yshift=-2mm, xshift=4mm},
  boxed title style={colback=forestgreen, arc=2pt},
  #1
}

% ── CUSTOM COMMANDS ───────────────────────────────────────────────────────────
\newcommand{\Hz}[1]{\texttt{\textbf{\textcolor{forestgreen}{#1}}}}
\newcommand{\tagbio}[0]{\colorbox{green!20}{\textcolor{green!50!black}{\small\bfseries BIOPHONY}}}
\newcommand{\taggeo}[0]{\colorbox{blue!15}{\textcolor{blue!70!black}{\small\bfseries GEOPHONY}}}
\newcommand{\taganthro}[0]{\colorbox{orange!20}{\textcolor{orange!60!black}{\small\bfseries ANTHROPOPHONY}}}

% ── TITLE ─────────────────────────────────────────────────────────────────────
\title{%
  {\color{darkgreen}\rule{\linewidth}{2pt}}\\[6pt]
  {\Huge\bfseries\color{forestgreen} Soundscapes \& Biodiversity}\\[4pt]
  {\large\color{darkgreen} Acoustic Environments, Animal Behaviour,\\
  Frequency Effects, and Restoration Strategies}\\[6pt]
  {\color{darkgreen}\rule{\linewidth}{2pt}}
}
\author{Research Synthesis --- Ecoacoustics \& Acoustic Ecology}
\date{May 2026}

% ══════════════════════════════════════════════════════════════════════════════
\begin{document}
% ══════════════════════════════════════════════════════════════════════════════

\maketitle
\thispagestyle{empty}

\vspace{0.5cm}

\begin{greenbox}[title=Abstract]
This document synthesises current scientific knowledge on soundscape ecology: how acoustic
environments are structured, how they influence living organisms across taxa, which frequency
ranges trigger which biological responses, and how acousticians and ecologists can leverage
sound as a tool to monitor and actively restore biodiversity. Research covered spans birds,
insects, fish, marine mammals, terrestrial mammals, amphibians, and plants, drawing on
peer-reviewed literature from 2003 to 2025.
\end{greenbox}

\vspace{0.5cm}
\tableofcontents
\newpage

% ══════════════════════════════════════════════════════════════════════════════
\section{What Is a Soundscape?}
% ══════════════════════════════════════════════════════════════════════════════

The term \textbf{soundscape} was coined by Canadian composer R.~Murray Schafer in the 1970s to
describe the total acoustic environment of a place --- the sum of all sounds, their
interactions, and the meaning those sounds carry for the beings that perceive them. In the
biological sciences, \textbf{soundscape ecology} was formalised as a discipline in 2011 with the
landmark article \textit{``Soundscape Ecology: The Science of Sound in the Landscape''}, which
defined it as the study of how sound from biological, geophysical, and anthropogenic sources can
illuminate the coupled dynamics of natural and human systems across space and time.

\subsection{The Three-Component Framework (Krause)}

The theoretical backbone of soundscape ecology was provided by ecologist and musician
\textbf{Bernie Krause}, who coined three foundational categories:

\vspace{6pt}
\begin{center}
\begin{tabularx}{0.96\linewidth}{>{\bfseries\color{forestgreen}}p{3cm} X}
\toprule
\rowcolor{rowshade}
Category & Definition \\
\midrule
\tagbio{} Biophony
  & The collective acoustic signatures produced by all living organisms in a habitat at a given
    moment: birdsong, frog choruses, insect stridulation, whale calls, bat ultrasound. The
    richest and most ecologically informative component of any soundscape. \\
\addlinespace
\taggeo{} Geophony
  & Abiotic natural sounds generated by physical forces: wind, flowing water, rain, thunder,
    ocean surf, and seismic activity. These form the evolutionary acoustic backdrop against which
    all biological communication evolved. \\
\addlinespace
\taganthro{} Anthropophony
  & All sounds generated by human activity, from music and speech to road traffic, aircraft,
    shipping, and industrial machinery. Anthropophony now dominates many of Earth's soundscapes,
    often to the detriment of biophony. \\
\bottomrule
\end{tabularx}
\end{center}

\subsection{The Acoustic Niche Hypothesis}

\begin{greenbox}[title=The Acoustic Niche Hypothesis --- Bernie Krause (1993)]
In a healthy, biodiverse habitat, species evolve to occupy \textbf{distinct acoustic niches} ---
partitioning the available frequency bandwidth and temporal windows so that their signals do not
mask one another. Each species adjusts its vocalisations by frequency and timing to minimise
interference with neighbouring acoustic territories. A rich, complex soundscape is therefore a
direct proxy for a biodiverse, well-functioning ecosystem. Conversely, when the soundscape
becomes impoverished or dominated by anthropophony, it signals ecological degradation --- often
years before visual biodiversity surveys register any change.
\end{greenbox}

\newpage

% ══════════════════════════════════════════════════════════════════════════════
\section{How Soundscapes Shape the Environment}
% ══════════════════════════════════════════════════════════════════════════════

Sound is not merely a \textit{byproduct} of life --- it is a fundamental ecological force.
The acoustic environment encodes continuous information about species presence, population health,
trophic interactions, seasonal phenology, water dynamics, and vegetation structure.
Disrupting it triggers cascades that ripple across multiple trophic levels.

\subsection{Key Figures}

\vspace{6pt}
\begin{center}
\begin{tabularx}{0.96\linewidth}{>{\centering\bfseries\color{forestgreen}\large}p{2.4cm} X}
\toprule
\rowcolor{rowshade}
Value & Finding \\
\midrule
$>$40\% & Of natural soundscapes on Earth have been so radically altered that many animal
          community members are presumed locally extinct \citep{krause2012}. \\
\addlinespace
30 dB   & Reduction in traffic noise achievable using green buffer zones combining dense
          vegetation with earthen berms. \\
\addlinespace
$\times$2 & Increase in fish abundance recorded on coral reef sites subjected to acoustic
             enrichment playback compared to control sites. \\
\addlinespace
$\times$7 & Higher coral larval settlement rates on degraded reefs exposed to healthy-reef
             sound playback via underwater speakers. \\
\bottomrule
\end{tabularx}
\end{center}

\subsection{The Cascade of Anthropogenic Noise}

When anthropophony intrudes into natural soundscapes, the consequences propagate through the
ecosystem via several interconnected pathways:

\begin{itemize}[leftmargin=1.5em, itemsep=4pt]
  \item \textbf{Communication masking:} human noise overlaps the frequency bands of animal vocal
        signals, preventing mates from locating each other, offspring from signalling parents,
        and conspecifics from coordinating group behaviour.

  \item \textbf{Habitat avoidance:} species abandon acoustically suitable areas due to noise
        alone, effectively contracting their range without any physical habitat change.

  \item \textbf{Trophic disruption:} predators lose acoustic detection ability; prey lose
        alarm-call networks; pollinator--plant relationships are impaired when buzz-pollination
        signals are masked.

  \item \textbf{Chronic physiological stress:} sustained noise elevates glucocorticoids
        (cortisol/corticosterone), suppresses immune function, reduces growth rates, and
        decreases reproductive success across taxa.

  \item \textbf{Acoustic niche collapse:} when low-frequency anthropophony saturates the
        soundscape, species vocalising in that range must either shift their signals --- at
        energetic cost --- or be excluded from the habitat entirely.
\end{itemize}

\begin{warnbox}
\textbf{Early-warning signal.} Long-term recordings by Bernie Krause spanning over five decades
show that acoustic degradation precedes visually detectable biodiversity decline by years. Passive
acoustic monitoring is therefore one of the most sensitive early-warning tools available to
conservation ecologists.
\end{warnbox}

\newpage

% ══════════════════════════════════════════════════════════════════════════════
\section{How Animals React to Soundscapes}
% ══════════════════════════════════════════════════════════════════════════════

Sound mediates nearly every critical behaviour in the animal kingdom --- mating, predator
detection, navigation, social cohesion, and habitat selection. Below is a taxon-by-taxon
survey of documented responses.

% ── BIRDS ─────────────────────────────────────────────────────────────────────
\subsection{Birds}

\begin{speciesbox}{Birds}
Noise pollution acts as a powerful \textbf{filter on bird communities}. Species with
low-frequency songs (typically larger-bodied birds) are selectively excluded from noisy
environments, while species capable of higher-frequency vocalisation persist. This restructures
entire communities based on vocal traits rather than ecological role.

\medskip
Urban birds display remarkable \textbf{vocal plasticity}. Great tits (\textit{Parus major}),
European blackbirds (\textit{Turdus merula}), and white-crowned sparrows
(\textit{Zonotrichia leucophrys}) have all been documented raising their minimum song frequency
in cities to avoid masking by low-frequency traffic noise. Additional documented adaptations
include:

\begin{itemize}[itemsep=2pt]
  \item Singing louder --- the \textit{Lombard effect} (amplitude increase in response to
        background noise);
  \item Shifting the timing of song bouts away from peak traffic hours;
  \item Shortening inter-bout intervals and simplifying song structure.
\end{itemize}

\medskip
These adaptations have reproductive costs. Species with narrow vocal plasticity or rigid
habitat requirements face population decline. A natural experiment at a closed international
airport showed that the dawn chorus recovered measurably within a single breeding season,
demonstrating both the sensitivity and the resilience of avian acoustics to noise pressure.
\end{speciesbox}

% ── INSECTS ───────────────────────────────────────────────────────────────────
\subsection{Insects}

\begin{speciesbox}{Insects}
Road-side grasshopper populations (\textit{Chorthippus biguttulus}) produce mating calls with
\textbf{elevated frequencies} that are less masked by vehicle noise. Crucially, this frequency
shift persists when individuals are transferred to a silent room, suggesting a genetic or
developmental fixation of the adaptation --- a rapid evolutionary response to chronic noise
exposure.

\medskip
Anthropogenic noise also disrupts insect--plant interactions:
\begin{itemize}[itemsep=2pt]
  \item The \textbf{chewing vibrations} of caterpillars trigger glucosinolate defence responses
        in plants; road noise at the same frequency range can suppress this detection.
  \item \textbf{Buzz-pollination} (sonication) relies on the resonant frequency of bee wing
        muscles ($\sim$200--400\,Hz) to release pollen; masking noise may impair this
        plant--pollinator coupling.
\end{itemize}
\end{speciesbox}

% ── FISH & MARINE LIFE ────────────────────────────────────────────────────────
\subsection{Fish and Marine Life}

\begin{speciesbox}{Fish \& Marine Life}
Healthy coral reefs generate rich soundscapes --- crackling snapping shrimp
(\textit{Alpheus} spp.), fish stridulation, parrotfish grazing. \textbf{Fish larvae use these
sounds as a navigational beacon} from open ocean to suitable reef habitat. Degraded reefs
produce impoverished soundscapes, directly reducing larval settlement and reef recovery.

\medskip
Anthropogenic underwater noise (shipping, seismic surveys, pile-driving, military sonar) causes:
\begin{itemize}[itemsep=2pt]
  \item Temporary and permanent threshold shifts in hearing sensitivity;
  \item Interrupted migration corridors and feeding routes;
  \item Strandings and spatial disorientation in cetaceans near sonar;
  \item Elevated cortisol and cytokine stress markers in bottlenose dolphins;
  \item Developmental abnormalities in larval zebrafish at sustained high-noise levels.
\end{itemize}

\medskip
Shipping noise specifically overlaps the low-frequency communication band of baleen whales,
shrinking their effective acoustic communication range by orders of magnitude --- from hundreds
of kilometres to just a few kilometres in heavily trafficked sea lanes.
\end{speciesbox}

% ── TERRESTRIAL MAMMALS ───────────────────────────────────────────────────────
\subsection{Terrestrial Mammals}

\begin{speciesbox}{Terrestrial Mammals}
Road noise impairs the ability of mongoose (\textit{Mungos mungo}) to detect predator-scent
cues --- a cross-modal cascade where acoustic masking degrades chemosensory vigilance and group
cohesion. Chronic noise exposure causes measurable cortisol elevation with downstream immune
suppression and reduced reproductive success.

\medskip
\textbf{Elephants} (\textit{Loxodonta africana}, \textit{Elephas maximus}) rely on infrasound
for long-distance herd coordination and threat communication across landscapes. They monitor
substrate-borne seismic signals through their feet and skeletal bones. Industrial low-frequency
and seismic anthropogenic noise directly overlaps and disrupts these channels.
\end{speciesbox}

% ── AMPHIBIANS ────────────────────────────────────────────────────────────────
\subsection{Amphibians}

\begin{speciesbox}{Amphibians}
Frog and toad choruses occupy specific temporal windows (dawn, dusk, rain events) and represent
key indicators of wetland health. Species richness in anuran communities is strongly correlated
with soundscape quality. Some species shift call frequency or timing in response to chronic noise
--- but these adjustments carry energetic costs and impair mate recognition. Frog community
composition has been shown to predict local noise levels with high accuracy, making anurans
effective bioacoustic sentinels.
\end{speciesbox}

\newpage

% ══════════════════════════════════════════════════════════════════════════════
\section{Frequencies: Which Hz Provokes Which Reaction?}
% ══════════════════════════════════════════════════════════════════════════════

The acoustic spectrum spans from seismic substrate vibration below 1\,Hz to ultrasound above
100\,kHz. Different frequency bands carry profoundly different biological significance.
Table~\ref{tab:frequencies} synthesises the key findings across taxa.

\begin{greenbox}
\textbf{Fundamental principle.} No frequency is inherently beneficial or harmful. Harm arises
from inappropriate \textit{amplitude} (too loud), inappropriate \textit{source} (masking natural
signals), or inappropriate \textit{timing} (overlapping with critical communication windows).
The goal of acoustic ecology is to restore balance across the full frequency landscape.
\end{greenbox}

\vspace{8pt}

\begin{longtable}{>{\Hz{}}p{2.7cm} p{2.8cm} p{8cm}}
\caption{Frequency ranges and their biological significance across taxa.}
\label{tab:frequencies} \\
\toprule
\textbf{Frequency} & \textbf{Classification} & \textbf{Key Users / Observed Effects} \\
\midrule
\endfirsthead
\toprule
\textbf{Frequency} & \textbf{Classification} & \textbf{Key Users / Observed Effects} \\
\midrule
\endhead
\bottomrule
\endfoot

$<$ 1 Hz
  & Seismic / substrate vibration
  & Elephants detect seismic rumbles through feet and bones over tens of km. Spiders and insects
    sense substrate vibrations for prey detection and mate signalling.
    Pile-driving and seismic surveys strongly disrupt these channels. \\
\rowcolor{rowshade}
1 -- 20 Hz
  & Infrasound
  & \textbf{Elephants}: coordinate herds over $>$10\,km; signal travels through air and soil
    simultaneously.
    \textbf{Baleen whales}: communicate across hundreds of km; blue whale calls reach
    14--30\,Hz at $>$180\,dB.
    \textbf{Rhinos, hippos, giraffes, okapis, peacocks, alligators}: conspecific
    long-distance communication.
    Biological tissue is particularly sensitive below $\sim$100\,Hz; industrial
    infrasound causes stress even when inaudible to humans. \\
20 -- 200 Hz
  & Low audible range
  & \textbf{Primary band of road traffic, shipping, and industrial noise} --- the dominant
    masker of biophony globally.
    Large-bodied birds (herons, owls, ravens) vocalise here and are selectively excluded
    from noisy environments.
    Most fish species hear in 20--1000\,Hz; shipping noise directly overlaps this range.
    Many frog species call in 100--500\,Hz. \\
\rowcolor{rowshade}
200 Hz -- 2 kHz
  & Mid audible range
  & Core vocal range of passerine birds; urban adaptation drives frequency minimum above
    $\sim$1.5--2\,kHz.
    Frog advertisement and territorial calls cluster here.
    \textbf{Plants}: most responsive in 100\,Hz--1\,kHz; 500\,Hz triggers auxin (IAA) and
    gibberellin (GA) biosynthesis in \textit{Arabidopsis}; 1\,kHz at 100\,dB enhances
    callus growth and drought tolerance in rice; 100\,Hz promotes root elongation. \\
2 -- 8 kHz
  & Upper audible range
  & Highest-information band for passerine communication; species retaining calls here in
    cities show better reproductive success.
    Grasshopper, cricket, and cicada stridulation; road-side populations have evolved
    elevated-frequency calls in this band.
    Snapping-shrimp reef ``crackling'' --- critical fish-larval recruitment cue ---
    peaks around 1--7\,kHz. \\
\rowcolor{rowshade}
8 -- 20 kHz
  & High audible / near-ultrasonic
  & Upper boundary of human hearing. Cats, dogs, rodents communicate here.
    Lower-frequency bat echolocation (free-tailed bats: 10--15\,kHz).
    Playback of reef sounds in this range combined with lower frequencies is used in
    acoustic reef restoration experiments. \\
20 -- 100 kHz
  & Ultrasound
  & \textbf{Bats}: echolocation from $\sim$20\,kHz (large species) to $>$100\,kHz (small
    species); high frequency reflects well off small objects enabling fine nocturnal
    navigation.
    \textbf{Toothed whales \& dolphins}: echolocation at 20--150\,kHz; boat noise can
    partially mask returning echoes.
    \textbf{Rodents}: ultrasonic social communication (30--90\,kHz) --- mouse courtship
    songs, rat pup distress calls.
    Some moths detect bat echolocation pulses and perform evasive manoeuvres. \\
$>$ 100 kHz
  & High ultrasound
  & Specialist bat species; some marine invertebrates. Largely beyond current anthropogenic
    interference. \\
\end{longtable}

\newpage

% ══════════════════════════════════════════════════════════════════════════════
\section{Plants in the Acoustic Landscape}
% ══════════════════════════════════════════════════════════════════════════════

Plants are not passive recipients of their acoustic environment. Research in
\textbf{plant bioacoustics} has established that sound waves induce mechanical vibrations in
plant tissues, alter fluid dynamics, increase membrane permeability, and trigger cascades of
gene expression.

\subsection{Mechanisms of Sound Perception in Plants}

Plants lack a centralised nervous system, but possess \textbf{mechanosensitive ion channels}
in cell membranes that convert pressure waves into electrochemical signals. At biologically
relevant intensities (50--100\,dB), sound causes measurable changes in:

\begin{itemize}[itemsep=3pt]
  \item Water and mineral nutrient uptake efficiency via increased membrane permeability;
  \item Stomatal aperture and transpiration rates;
  \item Expression profiles of genes governing primary metabolism and cell-wall modification
        (transcriptomic analyses confirm genome-wide responses);
  \item Biosynthesis and spatial distribution of phytohormones.
\end{itemize}

\subsection{Frequency-Specific Effects on Plants}

\begin{center}
\begin{tabularx}{0.96\linewidth}{>{\Hz{}}p{2.5cm} X}
\toprule
\rowcolor{rowshade}
Frequency & Documented Effect \\
\midrule
100 Hz
  & Promotes root elongation and cell division in root apical meristems of
    \textit{Arabidopsis thaliana} (15\,h/day exposure). \\
500 Hz
  & Induces production of auxin (IAA), gibberellin (GA), salicylic acid (SA), and jasmonic
    acid (JA) in \textit{Arabidopsis} --- simultaneous activation of growth and defence
    hormonal pathways. \\
\rowcolor{rowshade}
800 Hz -- 1.5 kHz
  & 1\,hour exposure increases drought stress tolerance in rice. \\
1 kHz @ 100 dB
  & Enhances callus growth in \textit{Dendranthema morifolium}; increases SOD
    (superoxide dismutase) antioxidant activity and soluble protein content in kiwi callus. \\
\rowcolor{rowshade}
20 Hz -- 1 kHz (general)
  & Broad range shown to enhance seed germination rates, seedling growth, and stress
    resistance across multiple crop species. \\
\bottomrule
\end{tabularx}
\end{center}

\subsection{Ecologically Meaningful Acoustic Signals for Plants}

\begin{itemize}[itemsep=4pt]
  \item \textbf{Bee buzz-pollination (sonication):} bees vibrate their flight muscles at
        $\sim$200--400\,Hz to release pollen from poricidal anthers. This is a finely tuned
        acoustic coupling between plant and pollinator; masking noise can disrupt it.

  \item \textbf{Herbivore chewing vibrations:} plants detect the characteristic vibration
        signature of caterpillar chewing and upregulate glucosinolate defences within minutes.
        Road noise overlapping these frequencies can suppress plant immune priming.

  \item \textbf{Rain and wind:} geophonic cues that trigger stomatal responses and
        phototropic / thigmotropic growth adjustments.
\end{itemize}

\begin{warnbox}
\textbf{Field caution.} Most positive results on sound-enhanced plant growth derive from
controlled laboratory settings. Applying artificial sound in complex field ecosystems risks
masking naturally occurring acoustic communication among species --- including the very signals
(bee buzz, caterpillar vibrations) that plants rely on for pollination and defence.
Any acoustic intervention must be carefully calibrated in context.
\end{warnbox}

\newpage

% ══════════════════════════════════════════════════════════════════════════════
\section{Monitoring Biodiversity Through Sound}
% ══════════════════════════════════════════════════════════════════════════════

\textbf{Ecoacoustics} has emerged over the past two decades as a cost-effective, scalable,
and non-invasive approach to biodiversity assessment. Autonomous recording units (ARUs) deployed
in the field capture continuous soundscape data that can be analysed with a growing toolkit
of acoustic indices and machine-learning species identification algorithms.

\subsection{Key Acoustic Indices}

A comprehensive 2025 practitioner's guide (\textit{Methods in Ecology and Evolution}) and a
systematic review of 187 studies (2003--2025) have formalised best practices. The most widely
used indices are:

\begin{itemize}[itemsep=4pt]
  \item \textbf{ACI --- Acoustic Complexity Index:} captures amplitude modulation irregularity
        across frequency bins. High ACI correlates with high biophony diversity.

  \item \textbf{ADI --- Acoustic Diversity Index:} measures evenness of sound energy across
        the spectrum; analogous to Shannon diversity applied to the acoustic domain.

  \item \textbf{AEI --- Acoustic Evenness Index:} inverse of ADI; high values indicate
        dominance by a single frequency band --- often a symptom of anthropophony.

  \item \textbf{BI --- Bioacoustic Index:} integrates spectral power in the 2--8\,kHz band,
        the core biophony window for birds and many insects.

  \item \textbf{NDSI --- Normalised Difference Soundscape Index:} ratio of biophony to
        anthropophony; a direct quantitative measure of acoustic naturalness.

  \item \textbf{Soundscape Saturation:} based on the acoustic niche hypothesis; reflects how
        fully acoustic niches are occupied --- a proxy for community species-packing.
\end{itemize}

\begin{bluebox}
\textbf{Scale of the field (2025).} A 2025 systematic review identified 187 ecoacoustic
biodiversity studies published between 2003 and May 2025. Passive acoustic monitoring (PAM)
was used in 89\% of cases. Bats were the most studied taxon, followed by birds, whole
soundscapes, amphibians, invertebrates, and mammals. The emerging frontier is
\textbf{soil ecoacoustics} --- detecting below-ground biodiversity through the acoustic
signatures of invertebrates and microorganisms in the substrate.
\end{bluebox}

\subsection{What Soundscapes Reveal}

\begin{itemize}[itemsep=3pt]
  \item Species richness and dominant taxa --- without visual surveys;
  \item Breeding phenology and seasonal dynamics across years;
  \item Canopy structure and vegetation complexity (via geophonic reverberation patterns);
  \item Rate of ecological recovery during and after restoration interventions;
  \item Level and spectral character of anthropogenic noise intrusion;
  \item Soil biodiversity --- an emerging application linking below-ground invertebrate
        acoustic activity to ecosystem function.
\end{itemize}

\newpage

% ══════════════════════════════════════════════════════════════════════════════
\section{Acoustic Solutions: Restoring and Growing Biodiversity}
% ══════════════════════════════════════════════════════════════════════════════

Acousticians, ecologists, and urban planners are developing a growing toolkit of interventions
--- from passive noise abatement to active acoustic seeding of life.

\subsection{Acoustic Enrichment and Playback}

Broadcasting recordings of healthy-ecosystem soundscapes in degraded areas attracts organisms
and accelerates recolonisation. This has been most rigorously demonstrated in marine ecosystems:

\begin{itemize}[itemsep=3pt]
  \item \textbf{Fish communities:} playback of healthy reef sounds on degraded coral habitat
        doubled overall fish abundance and increased species richness by 50\%, with
        improvements across all major trophic guilds (herbivores, planktivores, corallivores,
        piscivores) \citep{timber2019}.
  \item \textbf{Coral larvae:} underwater speakers playing healthy-reef soundscapes increased
        larval (\textit{Porites astreoides}) settlement rates up to 7$\times$ compared to
        silent controls \citep{vermeij2010}.
  \item \textbf{Terrestrial restoration:} broadcasting biophony in deforested areas attracts
        birds and bats, which passively transport seeds and fungal spores, accelerating
        vegetative recovery.
\end{itemize}

\subsection{Green Acoustic Buffers}

Dense vegetation absorbs and scatters sound waves. Combining tree belts, hedgerows, and
living walls with earthen berms (undulating topography) can reduce traffic noise by up to
30\,dB locally. Seoul's green transport corridors and London's ecological living-wall barriers
are leading urban examples. These buffers are \textbf{dual-benefit} interventions: they
attenuate noise \textit{and} create microhabitat for urban wildlife.

\subsection{Soundscape-Aware Urban Design}

Emerging ISO~12913 soundscape standards guide urban planners to introduce \textbf{positive
masking} natural sounds --- water features, wind-responsive vegetation, bird-attracting
plantings --- that shift the perceived acoustic environment from anthropophony-dominant to
biophony-dominant without necessarily requiring overall decibel reduction. This approach
simultaneously benefits human psychological wellbeing (reduced stress, improved restorativeness)
and urban biodiversity.

\subsection{Passive Acoustic Monitoring Networks}

Deploying arrays of autonomous recording units provides continuous, real-time ecosystem
surveillance. Combined with AI-powered species identification, this allows managers to track
biodiversity responses to interventions at landscape scale. Open platforms such as
\textit{ecoSound-web} democratise access to ecoacoustic analysis tools for citizen scientists
and conservation practitioners globally.

\subsection{Maritime Noise Reduction}

\begin{itemize}[itemsep=3pt]
  \item Reducing ship speed by 20\% cuts underwater radiated noise by up to 6\,dB, substantially
        expanding the acoustic communication range of whales.
  \item Propeller redesign and hull acoustic optimisation are engineering approaches under active
        discussion at the International Maritime Organization (IMO).
  \item \textbf{Quiet shipping corridors} during whale migration seasons are being trialled
        internationally as a regulatory tool.
\end{itemize}

\subsection{Acoustic Refuges and Protected Quiet Zones}

Designating areas free of motorised vehicles, aircraft, and industrial activity allows biophony
to recover undisturbed. The US National Park Service's \textit{Natural Sounds Program} and the
European Environmental Noise Directive both include mechanisms to protect acoustic habitat.
Research shows that even small quiet refuges (a few hectares) can sustain acoustic niche
diversity for dozens of species, provided they are embedded within a landscape matrix that
allows animal movement.

\subsection{Chronoacoustic Management}

Scheduling noisy human activities to avoid critical acoustic windows --- dawn chorus, dusk bat
emergence, nocturnal frog advertisement choruses, crepuscular mammal communication --- can
dramatically reduce biological cost without reducing the overall volume of activity. Time-based
noise restrictions represent one of the lowest-cost, highest-impact regulatory tools available
to land managers and urban authorities.

\subsection{Acoustic Agriculture}

Controlled exposure of crops to specific frequency ranges (500\,Hz -- 1\,kHz) in greenhouse or
precision-agriculture field settings can enhance germination, root development, and abiotic
stress tolerance, potentially reducing irrigation demand and agrochemical inputs. This must be
balanced against the risk of masking buzz-pollination and herbivore-detection pathways in
adjacent wild vegetation.

\begin{greenbox}[title=Emerging Consensus]
The most effective biodiversity restoration programmes combine \textbf{acoustic enrichment}
(active playback) with \textbf{physical habitat restoration} (revegetation, substrate
rehabilitation), \textbf{noise abatement} (barriers, quiet zones, temporal restrictions), and
\textbf{continuous passive acoustic monitoring} to track recovery in near-real time.
Sound is not merely a symptom of biodiversity --- it is a lever for restoring it.
\end{greenbox}

\newpage

% ══════════════════════════════════════════════════════════════════════════════
\section*{Conclusion}
\addcontentsline{toc}{section}{Conclusion}
% ══════════════════════════════════════════════════════════════════════════════

Soundscape ecology reveals that the acoustic environment is not a passive background to
biological life but an active ecological dimension --- one that organisms have co-evolved with
over millions of years, that encodes the health of ecosystems in real time, and that can be
degraded or restored through deliberate human choices.

The convergence of ecoacoustics, bioacoustics, urban design, and restoration ecology is opening
a new paradigm: managing landscapes not only for what they look like or what species lists they
contain, but for what they \textit{sound like}. A thriving soundscape --- rich in biophony,
structured by the acoustic niche hypothesis, and free from chronic anthropogenic masking --- is
both the signature and the engine of biodiversity.

For acousticians, the call to action is threefold:
\begin{enumerate}[itemsep=4pt]
  \item \textbf{Monitor}: deploy passive acoustic networks to provide baseline data and
        track ecosystem trajectories with unprecedented resolution.
  \item \textbf{Protect}: advocate for quiet zones, chronoacoustic restrictions, and
        noise-reduction standards in maritime and terrestrial policy.
  \item \textbf{Restore}: use acoustic enrichment as a scalable, cost-effective complement
        to conventional ecological restoration, anchored in rigorous monitoring to evaluate
        outcomes.
\end{enumerate}

% ══════════════════════════════════════════════════════════════════════════════
% BIBLIOGRAPHY
% ══════════════════════════════════════════════════════════════════════════════
\newpage
\section*{References}
\addcontentsline{toc}{section}{References}

\begin{thebibliography}{99}

\bibitem{krause2012}
Krause, B.~L. (2012). \textit{The Great Animal Orchestra}. Little, Brown \& Company.

\bibitem{pijanowski2011}
Pijanowski, B.~C., et al. (2011).
Soundscape ecology: The science of sound in the landscape.
\textit{BioScience}, 61(3), 203--216.

\bibitem{timber2019}
Timber, S., et al. (2019).
Acoustic enrichment can enhance fish community development on degraded coral reef habitat.
\textit{Nature Communications}, 10, 5414.
\url{https://www.nature.com/articles/s41467-019-13186-2}

\bibitem{vermeij2010}
Vermeij, M.~J.~A., et al. (2010).
Coral larvae move toward reef sounds.
\textit{PLOS ONE}, 5(5), e10660.

\bibitem{sager2007}
Slabbekoorn, H., \& Peet, M. (2003).
Birds sing at a higher pitch in urban noise.
\textit{Nature}, 424, 267.

\bibitem{noise_birds_pmc}
Francis, C.~D., et al. (2011).
Noise pollution filters bird communities based on vocal frequency.
\textit{PLOS ONE}, 6(11), e27052.
\url{https://pmc.ncbi.nlm.nih.gov/articles/PMC3212537/}

\bibitem{plants_pmc}
Rodrigues, J., et al. (2020).
Sound perception and its effects in plants and algae.
\textit{Plant Signaling \& Behavior}, 16(2), 1828674.
\url{https://pmc.ncbi.nlm.nih.gov/articles/PMC7671032/}

\bibitem{arabidopsis_roots}
Rodrigues, A., et al. (2021).
Sound waves promote \textit{Arabidopsis thaliana} root growth by regulating root phytohormone content.
\textit{PMC}.
\url{https://pmc.ncbi.nlm.nih.gov/articles/PMC8199107/}

\bibitem{symphonies2024}
Ghosh, S., et al. (2024).
Symphonies of growth: Unveiling the impact of sound waves on plant physiology and productivity.
\textit{PMC}.
\url{https://www.ncbi.nlm.nih.gov/pmc/articles/PMC11117645/}

\bibitem{acoustic_restoration}
Znidersic, E., et al. (2022).
Acoustic restoration: Using soundscapes to benchmark and fast-track recovery of ecological communities.
\textit{Ecology Letters}, 25(7), 1597--1607.
\url{https://onlinelibrary.wiley.com/doi/10.1111/ele.14015}

\bibitem{ecoacoustics2025}
Gitau, J., et al. (2025).
The role of ecoacoustics in monitoring ecosystem degradation and restoration.
\textit{Restoration Ecology}.
\url{https://onlinelibrary.wiley.com/doi/10.1111/rec.70168}

\bibitem{acoustic_index_guide}
Bradfer-Lawrence, T., et al. (2025).
The Acoustic Index User's Guide: A practical manual for defining, generating and understanding
current and future acoustic indices.
\textit{Methods in Ecology and Evolution}.
\url{https://besjournals.onlinelibrary.wiley.com/doi/full/10.1111/2041-210X.14357}

\bibitem{urban_ecology}
Urban soundscape monitoring for biodiversity. (2025).
\textit{Journal of Urban Ecology}, 11(1), juaf002.
\url{https://academic.oup.com/jue/article/11/1/juaf002/8088407}

\bibitem{royal_society}
Francis, C.~D., \& Barber, J.~R. (2020).
Direct and indirect effects of noise pollution alter biological communities in and near
noise-exposed environments.
\textit{Proceedings of the Royal Society B}, 287(1933).
\url{https://royalsocietypublishing.org/doi/10.1098/rspb.2020.0176}

\bibitem{dolphins_stress}
Kunc, H.~P., et al. (2021).
Anthropogenic sound exposure-induced stress in captive dolphins.
\textit{Frontiers in Marine Science}, 8, 606736.
\url{https://www.frontiersin.org/journals/marine-science/articles/10.3389/fmars.2021.606736/full}

\bibitem{soundscape_city}
Zhang, M., et al. (2022).
Soundscape in city and built environment: current developments and design potentials.
\textit{City and Built Environment}.
\url{https://pmc.ncbi.nlm.nih.gov/articles/PMC9833864/}

\bibitem{infrasound_elephants}
O'Connell-Rodwell, C.~E. (2007).
Keeping an "ear" to the ground: seismic communication in elephants.
\textit{Physiology}, 22(4), 287--294.

\bibitem{helsinki_animals}
University of Helsinki. (2023).
Animals communicate in frequencies we cannot hear.
\url{https://www.helsinki.fi/en/news/life-sciences/animals-communicate-frequencies-we-cannot-hear}

\end{thebibliography}

\end{document}
